\section{Local selection of low-frequency indices}

\begin{theorem}[Local choices in two frequency windows]
    \label{thm:two_window_low_frequency_selectors}
    \lean{Finset.exists_two_window_low_frequency_selectors}
    \leanok
    Let $I$ be a finite set and let $L,R\subseteq I$ be nonempty.
    There are functions
    $p_L\colon\N^L\to L$ and $p_R\colon\N^R\to R$ such that, for every
    frequency function $\nu\colon I\to\N$ and every $k\in\N$ with
    $\sum_{a\in I}\nu(a)\le k$,
    \begin{align}
        \nu\bigl(p_L(\nu|_L)\bigr)&\le\left\lfloor\frac{k}{|L|}\right\rfloor,
        \notag\\
        \nu\bigl(p_R(\nu|_R)\bigr)&\le\left\lfloor\frac{k}{|R|}\right\rfloor.
        \notag
    \end{align}
    The functions are chosen before $\nu$ and $k$, and each uses only the
    frequency function on its own window. The windows may overlap, and
    $k=0$ is allowed. This is the elementary selection step used in
    \cite[Lemma~4.2]{Arad2013AreaLaw} and in the simultaneous frequency
    argument of \cite[Proposition~5.7]{Arad2016RigorousRG}.
\end{theorem}
\begin{proof}
    \leanok
    For a nonempty window $S$, let $p_S(g)$ be a point where
    $g\colon S\to\N$ is minimal. For every global frequency function $\nu$,
    \begin{align}
        |S|\nu\bigl(p_S(\nu|_S)\bigr)
            &\le\sum_{a\in S}\nu(a)
            \le\sum_{a\in I}\nu(a)\le k.\notag
    \end{align}
    Since $|S|>0$, division and integrality give the claimed bound.
    Apply this argument separately to $L$ and $R$; no disjointness is used.
\end{proof}
